\documentclass[11pt]{article}
\usepackage{fontspec}
\usepackage[french,english]{babel}
\usepackage{amsmath}
\usepackage{amssymb}
\usepackage{amsthm}
\usepackage{mathrsfs}
\usepackage{tikz-cd}
\usepackage{enumitem}

\newcommand{\printedpage}[1]{\marginpar{\footnotesize p.~#1}}
\newcommand{\msp}[1]{\marginpar{\footnotesize [#1]}}
\newcommand{\circled}[1]{\tikz[baseline=(c.base)]\node[draw,circle,inner sep=1pt](c){#1};}
\newcommand{\Cat}{\mathrm{Cat}}
\newcommand{\Ob}{\mathrm{Ob}\,}
\newcommand{\hot}{\mathrm{hot}}

\begin{document}

% ===================== page 1 (PDF) =====================
\printedpage{[unnumbered]}
% head: (none printed)

\begin{center}
{\sffamily\LARGE THE DERIVATORS}

\bigskip

{\large\textbf{ALEXANDRE GROTHENDIECK}}

\vspace{3em}

{\large\textbf{Chapter XVI}}

\medskip

{\large\textbf{The basic localizers}}
\end{center}

\clearpage

% ===================== page 2 (PDF) =====================
\printedpage{[unnumbered]}
% head: (none printed)

This text was deciphered and transcribed in \LaTeXe{} by M. Künzer. It was edited by M. Künzer, J. Malgoire, and G. Maltsiniotis. The transcription is as faithful as possible to the manuscript. Only a few obvious minor corrections have been made. For the rare additions or comments by the editors, as well as for the original numbering of the manuscript pages, typewriter characters \texttt{[typewriter]} in brackets are used.

This edition is provisional. Remarks, comments and corrections are welcome. Send a message to:

\textbf{[e-mail address removed]}

\hfill G. Maltsiniotis\hspace*{2cm}

\clearpage

% ===================== page 3 (PDF) =====================
\printedpage{1}
% head: A. Grothendieck : Dérivateurs, Ch. XVI, Édition des Universités de Montpellier II et Paris VII
% foot: Version 18 mars 2019

\msp{page 1}

\section*{1\quad Axioms Loc(1) to Loc(3) (and variants): basic localizers}

I feel like taking up again the exposition of the axioms useful for a basic localizer
\[
W \subset \mathrm{Fl}(\Cat),
\]
guiding myself this time more by the geometric meaning of the axioms, than by the $\pm$ technical aliases one encounters (or does not encounter) in checking them. I am therefore revising in this spirit the presentation of XV §9.

\begin{description}
\item[Loc (1)] (\textit{saturation}).
\begin{itemize}
\item[\textit{a})] Iso $\subset W$.
\item[\textit{b})] If two among $f$, $g$, $gf$ are in $W$, the third too.
\item[\textit{c})] If $gf$ and $fg$ are in $W$, $f$ and $g$ too.
\end{itemize}
\end{description}

The conditions (\textit{a}) and (\textit{b}) alone, I am going to call them (for a general $W \subset \mathrm{Fl}\,\mathcal{M}$) \textit{weak saturation} of the localizer. The conditions (\textit{a}), (\textit{b}), (\textit{c}), \textit{mild saturation} (``mild saturation''). Other variants:

\begin{description}
\item[Loc (1 bis)] The conditions (\textit{a}), (\textit{b}), (\textit{c}) above, plus
\begin{itemize}
\item[\textit{d})] $W$ stable under direct factors.
\end{itemize}
\end{description}

This is of course implied by:

\begin{description}
\item[Loc (1 ter)] $W$ is \textit{strongly saturated, i.e.} if $f \in \Cat$ is such that $\hot_W(f)$ is an isomorphism, then $f \in W$.
\item[Loc (2)] (\textit{Final object}). If $X$ in Cat has a final object, then $X$ is $W$-aspheric, \textit{i.e.} $X \longrightarrow e$ is in $W$.
\end{description}

Sometimes, this axiom is a consequence of the following weakened variant.

\begin{description}
\item[Loc (2$_0$)] (Segment axiom). $\boldsymbol{\Delta}^1$ is $W$-aspheric.
\end{description}

\msp{page 2}

\begin{description}
\item[Loc (3)] (\textit{localization,} weak form). Let $f : X \longrightarrow S$ be an arrow in Cat which is $W$\textit{-aspheric, i.e.} such that $\forall\ s \in \Ob S$, the localized category $X/s$ over $S$ at $s$, $X/s = X \times_S S/s$, is $W$-aspheric. Then $f \in W$.\footnote[1]{In a final version, one should interchange the notations for Loc (3) and for Loc (3 bis), since the latter appears to me to be the more important of the two - it is satisfied for all the localizers that have been used (to my knowledge), and is valid for every localizer of the form $W_{\mathbf{D}}$, where $\mathbf{D}$ is a derivator (or just a prederivator satisfying Der (4,5) \texttt{[Der (2) is also needed; the numbering of the axioms is that of chapter I, p. 68]}).} % mark printed as superscript (1); [sic] utilisées -> utilisés, agreement with "localiseurs"
\end{description}

\textbf{Recall 1.} Loc (1, 2, 3) imply that we have
\[
W = W^\circ \qquad \textit{i.e. } f \in W \iff f^\circ \in W,
\]
\printedpage{2}
% head: A. Grothendieck : Dérivateurs, Ch. XVI, Édition des Universités de Montpellier II et Paris VII
% foot: Version 18 mars 2019
where we denote by $X \longmapsto X^\circ$, $f \longmapsto f^\circ$ the functor of passage to the opposite category, functor
\[
X \longmapsto X^\circ : \Cat \xrightarrow{\ \sim\ } \Cat
\]
which is an \textit{automorphism} of Cat (thus commutes with all $\varinjlim$ and $\varprojlim$ ). It follows that the set of axioms Loc (1, 2, 3) is equivalent to Loc (1, 2$^{\mathrm{opp}}$, 3$^{\mathrm{opp}}$), where Loc (2$^{\mathrm{opp}}$) is the ``dual'' axiom of Loc (2), namely that every object of Cat having an initial object is $W$-aspheric, and likewise Loc (3$^{\mathrm{opp}}$) the axiom that every $f : X \longrightarrow S$ which is $W$\textit{-coaspheric, i.e.} such that the colocalized categories $s\backslash X$ of $X$ over $S$ are $W$-aspheric, is in $W$.

\textbf{Definition 1.1.} By a \textit{basic localizer} we mean a subset $W$ of Fl(Cat) satisfying the axioms Loc (1, 2, 3) (\textit{i.e.} mild saturation, final object axiom, localization axiom).

\msp{page 3}

\textbf{Recall 1.2.} If $W$ is a basic localizer, then every \textit{homotopism} $f : X \longrightarrow Y$ of Cat is in $W$ (it is essentially for this reason that we imposed (\textit{c}) in mild saturation Loc (1)). If $f, g : X \rightrightarrows Y$ are homotopic (\textit{i.e.} belong to the same connected component of $\underline{\mathrm{Hom}}(X, Y)$), then $\hot_W(f) = \hot_W(g)$ and $f \in W \iff g \in W$.

Here is the strong form of Loc (3).

\begin{description}
\item[Loc (3 bis)] (\textit{Localization,} strong form). Let $S$ be in Cat and $f : X \longrightarrow Y$ a morphism in $\Cat/S$ such that ``$f$ is a $W$-equivalence locally over $S$'', \textit{i.e.} for every $s \in \Ob S$, the induced morphism
\[
f/s : X/s \longrightarrow Y/s
\]
for the localized categories of $X$, $Y$ over $S$ at $s$, is in $W$. Then $f \in W$.
\end{description}

The axiom Loc (3) is nothing but the special case of Loc (3 bis) corresponding to the case where $Y = S$.

\textbf{Proposition 1.3.} \textit{Let $W \subset \mathrm{Fl}(\mathcal{M})$. Then we have the following equivalences}
\[
\begin{aligned}
\text{Loc (1, 2, 3)} &\iff \text{Loc (1, $2_0$, $3'$)}\\
\text{Loc (1, 2, 3 bis)} &\iff \text{Loc (1, $2_0$, $3'$ bis),}
\end{aligned}
\]
\textit{where the axioms} Loc ($3'$) \textit{and} Loc ($3'$ bis) \textit{are formulated as follows:}

\msp{page 4}

\begin{description}
\item[Loc (3$'$)] \circled{\textit{a}} (Axiom of Cat-fibrations, weak direct form). Let $S$ be in Cat, $X$ a split fibered category over $S$, such that the fibers $X_s$ ($s \in \Ob S$) are $W$-aspheric. Then $X \longrightarrow S$ is in $W$.

\circled{\textit{b}} Let $X$ be in $\Cat/S$, then $i_X : X \longrightarrow \Phi_S(X)$ is in $W$\footnote[2]{(\textit{b}) concerns the \textit{comparison} of an arbitrary $X$ over $S$, with the split fibered category $\Phi_S(X)$ over $S$.}. % mark printed as superscript (2)
\end{description}

\clearpage

% ===================== page 5 (PDF) =====================
\printedpage{3}
% head: A. Grothendieck : Dérivateurs, Ch. XVI, Édition des Universités de Montpellier II et Paris VII
% foot: Version 18 mars 2019

\begin{description}
\item[Loc (3$'$ bis)] \circled{\textit{a}} (Axiom of Cat-fibrations, strong direct form). Let $S$ be in Cat, $X$, $Y$ two split fibered categories over $S$, $f : X \longrightarrow Y$ a cartesian $S$-functor compatible with the splittings, such that $\forall\ s \in \Ob S$, the induced functor on the fibers
\[
f_s : X_s \longrightarrow Y_s
\]
is in $W$. Then $f \in W$.

\circled{\textit{b}} As in Loc ($3'$) (\textit{b}).
\end{description}

\textbf{N.B.} In fact, assuming Loc (1, 2), the statements Loc ($3'$), Loc ($3'$ bis) imply the similar statements without splittings, and without even assuming $f$ cartesian in Loc ($3'$ bis), \textit{cf.} proof below.

\textsc{Proof.} \textit{a}) Loc (1, 2, 3) $\Longrightarrow$ Loc (1, $2_0$, $3'$). It suffices to prove Loc ($3'$). Let us prove Loc ($3'$ $a$). For this, under the conditions of Loc ($3'$ $a$), let us prove that $X \longrightarrow S$ is $W$-coaspheric, \textit{i.e.} that the categories
\[
s\backslash X = X \times_S s\backslash S
\]
are $W$-aspheric. It will indeed follow from this that $f \in W$ by virtue of Loc (3) applied to $f^\circ$ and of $W = W^\circ$. But
\msp{page 5}
let us consider the canonical inclusion
\[
i_s : X_s \lhook\joinrel\longrightarrow s\backslash X,
\]
it follows from the fact that $p : X \longrightarrow S$ is Cat-fibrant that this functor is $W$-aspheric (it even has a right adjoint, \textit{i.e.} the categories $X_s/u$, where $u \in \Ob(s\backslash X)$, have a final object). Thus by Loc (3) we have $i_s \in W$, and since $X_s$ is $W$-aspheric by hypothesis, likewise $s\backslash X$, q.e.d.

Let us prove Loc ($3'$ $b$): We note that $i_X : X \longrightarrow \Phi_S(X)$ is $W$\textit{-coaspheric} (it even has a left adjoint), hence $i_X^\circ$ is $W$-aspheric, hence in $W$, hence $i_X^\circ \in W^\circ$ since $W = W^\circ$. % [sic] i_X^\circ \in W^\circ -> i_X \in W (or i_X \in W^\circ), apparently intended: from i_X^\circ \in W one gets i_X \in W^\circ = W

$a$ bis) Loc (1, 2, 3 bis) $\Longrightarrow$ Loc (1, $2_0$, $3'$ bis). This time we have to prove only Loc ($3'$ bis). For (\textit{b}), this is already done. Remains to prove (\textit{a}). If $f$ is as in the statement Loc ($3'$ bis $a$), let us prove that for every $s$ in $S$, the induced morphism
\[
s\backslash f : s\backslash X \longrightarrow s\backslash Y
\]
is in $W$. Since $W = W^\circ$, we will therefore have $(s\backslash f)^\circ \in W$, and applying Loc (3 bis) to $f^\circ$, we find $f^\circ \in W$, hence $f \in W$, and we will be done. But $s\backslash f$ fits into the commutative square
\msp{page 6}
\[
\begin{tikzcd}
X_s \arrow[r, "f_s"] \arrow[d, "i_s^X"'] & Y_s \arrow[d, "i_s^Y"] \\
s\backslash X \arrow[r, "s\backslash f"] & s\backslash Y,
\end{tikzcd}
\]

\end{document}
